\documentclass[11pt]{article}

\usepackage[margin=0.62in]{geometry}
\usepackage{amsmath,amssymb}
\usepackage{enumitem}
\usepackage{xcolor}
\usepackage{tcolorbox}
\usepackage{fancyhdr}
\usepackage{lastpage}
\usepackage{hyperref}
\usepackage{microtype}
\usepackage{array}

\definecolor{uwpurple}{HTML}{4B2E83}
\definecolor{uwgold}{HTML}{B7A57A}
\definecolor{softpurple}{HTML}{F4F1FA}
\definecolor{softgold}{HTML}{FBF7EA}
\definecolor{deepgold}{HTML}{8A6F2A}

\tcbuselibrary{breakable,skins}

\hypersetup{
    colorlinks=true,
    linkcolor=uwpurple,
    urlcolor=uwpurple
}

\pagestyle{fancy}
\fancyhf{}
\lhead{\textbf{MATH 394: Probability I}}
\rhead{\textbf{Midterm Examination}}
\cfoot{\thepage\ of \pageref{LastPage}}
\setlength{\headheight}{14pt}

\setlength{\parindent}{0pt}
\setlength{\parskip}{0.4em}

\newcommand{\course}{MATH 394: Probability I}
\newcommand{\instructor}{Arman Jahangiri}
\newcommand{\department}{Department of Mathematics}
\newcommand{\university}{University of Washington}

\newtcolorbox{problemheader}{
    colback=softpurple,
    colframe=uwpurple,
    boxrule=0.8pt,
    arc=2mm,
    left=2mm,
    right=2mm,
    top=1mm,
    bottom=1mm
}

\newtcolorbox{policybox}{
    breakable,
    colback=softgold,
    colframe=deepgold,
    boxrule=0.8pt,
    arc=2mm,
    left=2.5mm,
    right=2.5mm,
    top=1.5mm,
    bottom=1.5mm
}

\newcommand{\workspace}[1]{%
    \vspace{0.4cm}
    \begin{tcolorbox}[
        height=#1,
        colback=white,
        colframe=uwgold,
        boxrule=0.5pt,
        arc=1mm,
        title={\small\textbf{Work space}},
        coltitle=uwpurple,
        colbacktitle=softpurple
    ]
    \end{tcolorbox}
}

\begin{document}

\begin{titlepage}
\centering
\vspace*{0.05cm}

{\Large \textbf{\university}}\\
{\large \department}

\vspace{0.28cm}

{\LARGE \color{uwpurple}\textbf{Midterm Examination}}\\[0.15cm]
{\Large \textbf{\course}}\\[0.15cm]
{\large Summer 2026 \quad | \quad Instructor: \instructor}

\vspace{0.22cm}

\begin{tcolorbox}[
    width=0.94\textwidth,
    colback=softpurple,
    colframe=uwpurple,
    boxrule=0.9pt,
    arc=3mm
]
\centering
\renewcommand{\arraystretch}{1.15}
\begin{tabular}{>{\bfseries}p{0.19\textwidth}p{0.68\textwidth}}
Date: & July 22, 2026 \\
Answering time: & 60 minutes \\
Total points: & 60 \\
\end{tabular}
\end{tcolorbox}

\vspace{0.18cm}

\begin{tcolorbox}[
    width=0.94\textwidth,
    colback=white,
    colframe=uwgold,
    boxrule=0.8pt,
    arc=3mm
]
\renewcommand{\arraystretch}{1.35}
\begin{tabular}{>{\bfseries}p{0.22\textwidth}p{0.67\textwidth}}
Name: & \hrulefill \\
Student ID: & \hrulefill \\
Signature: & \hrulefill \\
\end{tabular}
\end{tcolorbox}

\vspace{0.15cm}

\begin{policybox}
\small
\textbf{Instructions}
\begin{itemize}[leftmargin=*,itemsep=0.05em,topsep=0.15em,parsep=0pt]
    \item This examination contains \textbf{five problems}. Verify that your copy is complete.
    \item Show all work and justify your answers. Unsupported answers may receive little or no credit.
    \item Clearly define any events or random variables that you introduce.    
    \item Do not remove pages from the examination booklet.
\end{itemize}

\medskip
\textbf{Permitted materials}

Examinations are closed-book and closed-notes. Students may use:
\begin{itemize}[leftmargin=*,itemsep=0.05em,topsep=0.12em,parsep=0pt]
    \item one handwritten 8.5 $\times$ 11 inch reference sheet (front and back), and
    \item a non-programmable non-graphing calculator capable of standard exponential and logarithmic computations (e.g. TI-30X IIS).
\end{itemize}
No other electronic devices, calculators, communication devices, or external resources may be used during examinations.
\end{policybox}

\vspace{0.10cm}

\renewcommand{\arraystretch}{1.18}
\begin{tabular}{|>{\centering\arraybackslash}p{0.16\textwidth}|>{\centering\arraybackslash}p{0.16\textwidth}|}
\hline
\textbf{Problem} & \textbf{Score} \\
\hline
1 & /10 \\
\hline
2 & /10 \\
\hline
3 & /10 \\
\hline
4 & /15 \\
\hline
5 & /15 \\
\hline
\textbf{Total} & \textbf{/60} \\
\hline
\end{tabular}

\vspace{0.08cm}
{\footnotesize By signing above, you affirm that the work is your own and that you followed the examination rules.}
\end{titlepage}








% ------------------------------------------------------------
\begin{problemheader}
\textbf{Problem 1. } \hfill \textbf{5 points}
\end{problemheader}
A student organization has 7 mathematics graduate students, 5 statistics graduate students, and 6 undergraduate students. A committee of 5 students is selected uniformly at random from all committees of 5.

Find the probability that the committee contains exactly 2 mathematics graduate students, exactly 1 statistics graduate student, and exactly 2 undergraduate students.
\hrule
 

\newpage

% ------------------------------------------------------------
\begin{problemheader}
\textbf{Problem 2.} \hfill \textbf{10 points}
\end{problemheader}

A spam filter is designed by looking at commonly occurring phrases in spam. Suppose that $80\%$ of email is spam. In $10\%$ of the spam emails, the phrase ``free money'' is used, whereas this phrase is only used in $1\%$ of non-spam emails. A new email has just arrived, which does mention ``free money.'' What is the probability that it is spam?
 
\hrule
 

\newpage
% ------------------------------------------------------------
\begin{problemheader}
\textbf{Problem 3. } \hfill \textbf{10 points}
\end{problemheader}

Let $A$ and $B$ be events.

\begin{enumerate}[label=(\alph*),leftmargin=*]
    \item Let $S$ be a sample space and $A_1,\dots, A_n\in  \mathcal{F}$ be events. Explain what it means for $A_i's$ to form a partition of $S$,  \hfill \textbf{(3 points)}
    \item Suppose $P(A) = 0.6$ and $P(B)=0.5$. Can $A$ and $B$ be disjoint? \hfill \textbf{(3 points)}
   
    \item Suppose that $A$ is independent of itself. Use the definition of independence to conclude that $P(A)=0$ or $P(A)=1$. \hfill \textbf{(4 points)}
\end{enumerate}
 
\hrule
 
\newpage
% ------------------------------------------------------------
\begin{problemheader}
\textbf{Problem 4. } \hfill \textbf{15 points}
\end{problemheader}

Suppose that $X$ has probability density function
\[
f_X(x)=
\begin{cases}
 c\,x(2-x), & 0<x<2,\\
 0, & \text{otherwise},
\end{cases}
\]
where $c$ is a constant.

\begin{enumerate}[label=(\alph*),leftmargin=*,itemsep=0.7em]
    \item Find $c$. \hfill \textbf{(3 points)}
    \item Find the cumulative distribution function $F_X(x)$ for all real $x$. \hfill \textbf{(5 points)}
    \item Find $P(X>3/2)$. \hfill \textbf{(2 points)}
    \item Find $\mathbb{E}[X]$ and $\operatorname{Var}(X)$. \hfill \textbf{(5 points)}
\end{enumerate}

\hrule

\newpage
 




\begin{problemheader}
\textbf{Problem 5.} \hfill \textbf{20 Points}
\end{problemheader}

A company has found that \(40\%\) of customers who visit its website make a purchase. Suppose that \(200\) customers visit the website independently on a particular day.

Let \(X\) denote the number of customers who make a purchase.

\begin{enumerate}[]

\item Determine whether a \textbf{Poisson approximation} or a
\textbf{normal approximation} is more appropriate by verifying the
required conditions. Then clearly state the approximating distribution;
that is, provide \(\lambda\) if a Poisson approximation is appropriate,
or provide \(\mu\) and \(\sigma^2\) if a normal approximation is
appropriate.   \hfill \textbf{(5 points)}

\item  Use the approximation from part (a) to estimate each of the following:
\begin{enumerate}[]
    \item  \(\mathbb{P}(70<X<90)\)  \hfill \textbf{(5 points)}
    \item  \(\mathbb{P}(X>90)\)  \hfill \textbf{(5 points)}
    \item  \(\mathbb{P}(X=80)\)  \hfill \textbf{(5 points)}
\end{enumerate}

\end{enumerate}

\hrule

  

\end{document}
